\begin{quote}\small
\textbf{Аннотация.} Токены облачных моделей дороги: в агентской работе каждый вызов заново оплачивает чтение всего
накопленного контекста, и у медианного поставщика API на это уходит 85\% цены вызова. Открытую модель уровня
Qwen3.8-27B пользователь может запустить на своём компьютере, но с ростом контекста скорость работы падает,
потому что каждую новую порцию промпта она читает поверх всего накопленного контекста. На M3 Max скорость обработки
промпта (prefill) падает со 174 до 28 токенов в секунду при росте окна контекста с 8 тысяч до 181 тысячи токенов, и
даже порция в 10 тысяч новых токенов контекста вызывает задержку в 4--6 минут,
делая работу в реальном времени невозможной. Мы
переносим prefill на сервер: сервер держит копию сессии, читает только новые токены и присылает клиенту получившиеся
KV-кэш и рекуррентное состояние, а генерация остаётся на компьютере пользователя. Это быстро: по каналу в 1~Гбит/с
порция в 10 тысяч токенов на глубине 200 тысяч доходит до клиента за 5 секунд вместо 6.6 минуты. Это дёшево: на
журналах 1.5 миллиона вызовов агентов гибрид в 2.8--4.6 раза дешевле медианного поставщика API той же модели, хотя и
в 1.6--2.6 раза дороже самого дешёвого (в расчёте пользователь отдаёт открытой модели каждый пятый запрос, около 500 в день, а остальные отправляет более
сильным облачным моделям; Mac с чипом Max или Ultra у него уже есть).
Одной RTX PRO 6000 на сервере компании хватит на 64 активно работающих сотрудника, если открытая модель обрабатывает
пятую часть их запросов, и на 12, если всю работу.
Скорости клиента измерены, сервер и экономика рассчитаны воспроизводимым движком, обмен памятью проверен на прототипе с
малой моделью; качество генерации мы не оценивали.
\end{quote}
